\documentclass[11pt]{article}

\usepackage[a4paper,margin=1in]{geometry}
\usepackage{amsmath,amssymb,amsthm,booktabs}
\usepackage{microtype}

\newtheorem{theorem}{Theorem}
\newtheorem{lemma}[theorem]{Lemma}

\title{Exact Target-Sign Partial Ranges for\\
Atlas 153, 171, 174, and 188}
\author{}
\date{}

\begin{document}

\maketitle

\section{Statement}

For a finite graph $H$ and a graphon $W$, write
\[
 p=t(K_2,W),
 \qquad
 \Phi_H(p)
 =(1-p)^{v(H)}\chi_H\!\left(\frac1{1-p}\right).
\]
The catalogue asks whether
\[
 t(H,W)\geq\Phi_H(p)
\]
throughout the Tur\'an range.  For the four graphs in this note the
chromatic target remains nonpositive for a nontrivial interval beyond
$p=7/12$.  Nonnegativity of homomorphism densities therefore settles that
whole interval without any structural assumption on $W$.

\begin{theorem}[Exact nonpositive-target ranges]\label{thm:main}
For $H$ equal to Atlas $153,171,174$, or $188$, let $Q_H$ be as follows:
\[
\begin{array}{c|c|c}
H&Q_H(p)&\text{unique real zero }\alpha_H\\ \hline
153&7p^3-12p^2+8p-2&0.6068701814<\alpha_{153}<0.6068701815\\
171&11p^3-20p^2+13p-3&0.6271894887<\alpha_{171}<0.6271894888\\
174&13p^3-25p^2+17p-4&0.5924019902<\alpha_{174}<0.5924019903\\
188&17p^3-33p^2+22p-5&0.6108707597<\alpha_{188}<0.6108707598
\end{array}
\]
Then every graphon $W$ of edge density
\[
 \frac12\leq p\leq\alpha_H
\]
satisfies
\[
 t(H,W)\geq\Phi_H(p).
\]
In particular, the rounded-down decimal endpoints displayed in the table
give certified closed subintervals.
\end{theorem}

\section{Chromatic targets}

The graph6 codes and chromatic polynomials are
\[
\begin{array}{c|c|c}
H&\text{graph6}&\chi_H(x)\\ \hline
153&\texttt{Ehf\_}&x(x-1)(x-2)(x^3-5x^2+9x-7)\\
171&\texttt{Ehew}&x(x-1)(x-2)(x^3-6x^2+13x-11)\\
174&\texttt{EtTg}&x(x-1)(x-2)(x^3-6x^2+14x-13)\\
188&\texttt{EzNG}&x(x-1)(x-2)(x^3-7x^2+18x-17).
\end{array}
\]
Substituting $x=1/(1-p)$ and multiplying by $(1-p)^6$ gives, in every
case,
\begin{equation}\label{eq:target}
 \Phi_H(p)=p(2p-1)Q_H(p),
\end{equation}
with the four polynomials $Q_H$ from Theorem~\ref{thm:main}.

\section{The sign certificate}

\begin{lemma}\label{lem:sign}
Each $Q_H$ has exactly one real zero, this zero is larger than $1/2$, and
\[
 Q_H(p)\leq0
 \qquad\left(\frac12\leq p\leq\alpha_H\right).
\]
\end{lemma}

\begin{proof}
The cubic discriminants, in Atlas order $153,171,174,188$, are
\[
 -44,\qquad -31,\qquad -59,\qquad -23.
\]
They are strictly negative, so every $Q_H$ has exactly one real root.
Direct substitution gives the common value
\[
 Q_H\!\left(\frac12\right)=-\frac18.
\]
Every $Q_H$ has positive leading coefficient.  Its unique real zero is
therefore larger than $1/2$, and the polynomial is negative to the left of
that zero.

For completeness, the displayed decimal brackets are exact rational sign
checks: substituting their left endpoints gives respectively
\[
\begin{split}
 &-\frac{7725129772420851499}{125000000000000000000000000000},\\
 &-\frac{6439081474941456867}{1000000000000000000000000000000},\\
 &-\frac{330401596447139437}{125000000000000000000000000000},\\
 &-\frac{35667743563068492059}{1000000000000000000000000000000},
\end{split}
\]
while substituting the right endpoints gives respectively
\[
\begin{split}
 &\frac{440980018562085029}{8000000000000000000000000000},\\
 &\frac{161939305390749191}{1953125000000000000000000000},\\
 &\frac{104013296381851155251}{1000000000000000000000000000000},\\
 &\frac{4464622023299603783}{125000000000000000000000000000}.
\end{split}
\]
Thus the signs at every bracket endpoint are certified in exact rational
arithmetic.
\end{proof}

\begin{proof}[Proof of Theorem~\ref{thm:main}]
Let $1/2\leq p\leq\alpha_H$.  Then
\[
 p\geq0,
 \qquad
 2p-1\geq0,
 \qquad
 Q_H(p)\leq0
\]
by Lemma~\ref{lem:sign}.  Equation~\eqref{eq:target} therefore gives
$\Phi_H(p)\leq0$.  On the other hand, the integrand defining a graph
homomorphism density is nonnegative, so $t(H,W)\geq0$.  Hence
\[
 t(H,W)\geq0\geq\Phi_H(p),
\]
as required.
\end{proof}

\section{Scope}

The argument is valid for every graphon and uses no regularity,
finite-step, topology, or perturbative assumption.  It does not settle any
density above $\alpha_H$, where the target is positive.  Accordingly these
four catalogue rows are partial rather than positive.

\end{document}
